\chapter{Endliche Einbettungen}
\label{chap:b37-guide-representations}

\section{Einbettungen mit vorgeschriebenen Erzeugern}

\begin{definition}[Einbettung, Erzeugnis und zwei Erzeuger]
\label{def:b37-guide-generators}
Eine \emph{Einbettung} von Halbgruppen ist ein injektiver
Halbgruppenhomomorphismus.

Für eine Teilmenge $X$ einer Halbgruppe $A$ ist
\[
 \langle X\rangle_A
 =\{x_1\cdots x_m\mid m\geq1,\ x_1,\ldots,x_m\in X\}
\]
das \emph{Erzeugnis} von $X$. Für $X\ne\varnothing$ ist dies die von
$X$ erzeugte Unterhalbgruppe. Insbesondere ist
$\langle\varnothing\rangle_A=\varnothing$.
Eine Halbgruppe $A$ heißt \emph{von zwei Elementen erzeugt}, wenn
$A=\langle\{p,q\}\rangle_A$ für geeignete $p,q\in A$ gilt.
\end{definition}

\begin{definition}[Relative Freiheit auf zwei Erzeugern]
\label{def:b37-guide-relative-free}
Eine Halbgruppe $F=\langle\{p,q\}\rangle_F$ heißt
\emph{relativ frei auf den Erzeugern $p,q$}, wenn für alle $a,b\in F$
genau ein Halbgruppenendomorphismus $\psi:F\to F$ mit
\[
 \psi(p)=a,\qquad \psi(q)=b
\]
existiert. Ein Halbgruppenendomorphismus ist dabei ein
Halbgruppenhomomorphismus der Halbgruppe in sich selbst.
\end{definition}

\begin{theorem}[Endliche Einbettungen mit wenigen Erzeugern]
\label{thm:b37-guide-finite-embeddings}
Zu jeder endlichen Halbgruppe $S$ existieren die folgenden Einbettungen:
\begin{enumerate}
 \item Eine Einbettung $S\hookrightarrow R$ in eine nichtleere endliche
 Halbgruppe $R$, in der jedes Element regulär ist und die von zwei
 Elementen erzeugt wird. \cite{Higgins2017}
 \item Eine Einbettung $S\hookrightarrow F$ in eine nichtleere endliche
 Halbgruppe $F$, die relativ frei auf zwei Erzeugern
 ist. \cite{Bergman2004}
 \item Eine Einbettung $S\hookrightarrow T$ in eine nichtleere endliche
 Halbgruppe $T$ mit höchstens drei idempotenten Erzeugern:
 \[
 T=\langle\{e_1,e_2,e_3\}\rangle_T,
 \qquad e_i^2=e_i\quad(i=1,2,3).
 \]
 \cite[S.~62]{Higgins2017}
\end{enumerate}
\end{theorem}
